% ============================================================
%  MÓDULO 24 — Kit de aula: mapa mental e slides (R101)
%  Arquivos em material-didatico, gerados no mesmo caderno
% ============================================================
\chapter{Kit de Aula: Mapa Mental e Slides}

% R901-PONTE-ABERTURA
As pranchas ofereceram modelos de representação. Agora o leitor pode reorganizar o mesmo caso num mapa mental ou numa apresentação, mantendo fontes e limites ao reduzir o texto.


\roteirodidatico
{Você já tem as seis figuras do atlas; agora ganha o mapa que liga tudo e os slides para projetar em sala. Este capítulo curto apresenta os dois arquivos, diz onde estão e como usá-los sem improviso.}
{\emph{Mapa mental} é a árvore de ideias em JSON; \emph{Slide-deck} é a apresentação de 10 páginas em PDF e PPTX editável.}
{Abra o JSON para ver a árvore, projete o PDF em sala e edite o PPTX se precisar adaptar.}
{Arquivos gerados por IA a partir dos roteiros curtos: o texto principal foi conferido, mas a aula continua sendo sua — revise antes de projetar.}
{Que ramo do mapa você usaria para abrir uma aula de 10 minutos?}

\casoguiado{Carla projeta, Bruno edita}
{Carla abre o PDF e projeta a capa com o prédio de 8 andares. Bruno abre o PPTX, troca um exemplo pelo da sua turma e salva como nova versão com data. Os dois usam o mesmo kit: ela projeta o pronto, ele adapta o editável. Nenhum dos dois cita número do slide sem conferir o capítulo correspondente.}

\section{O mapa mental em 6 ramos}

\elemento{Mapa em JSON}{\metainfo{ID}{KIT-MAP} \hfill \metainfo{Arquivo}{\pth{material-didatico/mindmap-livro-core.json}}}

\begin{descricao}
Raiz \textbf{Sistema de Trabalho e Qualidade} com 6 ramos: (1) Estrutura, a casa de 8 andares; (2) Ciclo de Pedidos, as 4 fases com o rastro da entrega; (3) Fluxo da Encomenda, os 7 passos; (4) Gestão de Tarefas e Registro, lousa, caderno e verificação real (código, estado, resultado); (5) Controle de Qualidade, combinado antes, prova depois, fiscal fechado; (6) Métricas e Dados, 62 por cento na área difícil, 94 na fácil, com rastreabilidade. O JSON abre em qualquer editor de texto e no visualizador do NotebookLM (caderno \pth{419f77cb-6233-44a7-9ba3-b0ba19219135}).
\end{descricao}

\begin{funcao}
Dar a visão de 1 minuto do livro inteiro: do todo aos ramos, dos ramos aos capítulos, sem pular a verificação.
\end{funcao}

\section{Os slides em 10 páginas}

\elemento{Apresentação pronta e editável}{\metainfo{ID}{KIT-SLD} \hfill \metainfo{Arquivos}{\pth{slides-livro-core.pdf} e \pth{slides-livro-core.pptx}}}

\begin{descricao}
Abertura com o prédio (capa: Arquitetura, Ciclos e Qualidade), quadro de parâmetros (8 andares, 4 fases, 7 passos, 3 travas, 12 passos) e páginas por tema seguindo os roteiros curtos. Texto principal conferido visualmente em 11/10/2026 (acentos e números corretos nas páginas amostradas). Use o PDF para projetar e o PPTX para adaptar à sua turma, sempre com data e versão na primeira página editada.
\end{descricao}

\begin{aviso}
Limite do kit. Slides e mapa resumem; não substituem capítulos, arquivos nem testes. Erro de digitação residual em slide é possível: antes de projetar, folheie as 10 páginas uma vez e corrija no PPTX. Em divergência, valem capítulo e fonte primária.
\end{aviso}

\begin{resultado}
\textbf{Exercícios.} (1) Abra o JSON e conte os 6 ramos sem olhar a lista acima. (2) Folheie as 10 páginas do PDF e anote uma por capítulo do livro. (3) No PPTX, adapte um exemplo à sua realidade e salve com nova data.
\end{resultado}
\fontecore{livro-core/material-didatico}{caderno-419f77cb-mindmap-slides}

\section[Leitura guiada]{Leitura guiada: adaptar a aula preservando a proveniência}
\begin{descricao}
\textbf{O fenômeno.} Um slide e um mapa condensam conteúdo. Ao adaptar um exemplo ou número, muda-se o alcance do material e pode-se perder sua fonte original.
\end{descricao}
\begin{definicao}
\textbf{Grandezas e definições.} O mapa em JSON representa relações editoriais; PDF é a versão para projeção; PPTX é a fonte editável da apresentação. Data e identificação da versão distinguem a adaptação do kit recebido.
\end{definicao}
\begin{metodo}
\textbf{Dedução passo a passo.} Confira os arquivos; leia os ramos e slides; relacione-os aos capítulos; selecione o exemplo da aula; altere a cópia editável; conserve fontes e limites; exporte e reveja a projeção.
\end{metodo}
\begin{resultado}
\textbf{Exemplo resolvido.} Carla adapta um slide de métricas. Antes de projetar, consulta a medição original e inclui data, universo e limitações. Se não encontra esses dados, apresenta o tema como questão a investigar, sem converter um número do desenho em resultado atual.
\end{resultado}
\begin{resultado}
\textbf{Ordem de grandeza e verificação.} Os seis ramos e dez páginas são quantidades declaradas pelo kit original, a conferir nos arquivos. Uma adaptação com outra quantidade recebe novo registro; não se força a contagem anterior.
\end{resultado}
\begin{aviso}
\textbf{Limite de validade.} O kit é apoio didático. Não substitui a leitura, a execução dos testes nem a revisão científica. A inspeção atual do livro não implica nova conferência de todas as páginas do PPTX externo.
\end{aviso}
\begin{resultado}
\textbf{Exercícios.} Adapte um exemplo e indique sua fonte; registre a versão; confira a página exportada; explique uma alegação que a apresentação deliberadamente não demonstra.
\end{resultado}

% R901-PONTE-FECHO
Reduzir uma explicação exige saber o que a sustenta. Antes de concluir a reprodução, voltaremos às teorias citadas para distinguir inspiração conceitual, implementação e resultado observado.
